\question{25}{
Una tienda debe despachar un pedido de \emph{exactamente} $P \ge 0$ unidades de un producto. En bodega hay $n$ cajas cerradas: la caja $i$ contiene $a_i \ge 1$ unidades y despacharla cuesta $c_i \ge 0$ pesos. Cada caja se despacha a lo sumo una vez y no se puede abrir para sacar unidades sueltas. Determine el costo mínimo para despachar exactamente $P$ unidades, o determine que no es posible.

E.g., con $n=4$, $a = (3, 5, 4, 6)$ y $c = (5, 6, 4, 9)$:
\begin{itemize}
  \item $P=0$: $0$ (no se despacha ninguna caja).
  \item $P=2$: No es posible.
  \item $P=7$: $9$ (cajas $1$ y $3$).
  \item $P=10$: $13$ (cajas $3$ y $4$).
  \item $P=11$: $15$ (cajas $2$ y $4$).
  \item $P=13$: $18$ (cajas $1$, $3$ y $4$).
\end{itemize}
}

\begin{enumerate}
  \item (10 décimas) Defina una relación de recurrencia que resuelva el problema. Explique en palabras qué significa la función y qué significa cada variable. Sea sucinto.
  \item (5 décimas) Dibuje y explique el grafo de dependencias. Determine cuál será la complejidad espacial de su solución.
  \item (10 décimas) Implemente una solución de programación dinámica, ya sea bottom-up o top-down, en Python, Java, C, C++, JavaScript o TypeScript. Determine la complejidad temporal de su solución.
\end{enumerate}

\bigskip
